\chapter{Các thành phần của hệ thống}
\label{chap:comp}

\section{\texorpdfstring{Deep Learning $\Rightarrow$ Agent}{Deep Learning => Agent}}
\label{sec:dl}

\subsection{Dữ liệu và bài toán}

Khối ``DB có sẵn'' là bộ dữ liệu NHANES 2011–2018 của CDC \cite{nhanes}: \textbf{20.646 người trưởng thành}, trong đó 19,8\% mắc ĐTĐ (đã được chẩn đoán, hoặc có FPG ≥ 126 mg/dL, hoặc HbA1c ≥ 6,5\%). Dữ liệu được nạp vào bảng \texttt{nhanes\_samples} của SQLite, và mô hình đọc lại từ đây khi huấn luyện, đúng tinh thần ``DB → Deep learning'' của sơ đồ. Mô hình chỉ dùng \textbf{13 yếu tố nguy cơ}: tuổi, giới, BMI, vòng eo, huyết áp tâm thu và tâm trương, tăng huyết áp, rối loạn mỡ máu, tiền sử gia đình, hút thuốc, vận động, bệnh tim mạch, đột quỵ. Mô hình \emph{không} nhận glucose hay HbA1c. Đây là lựa chọn có chủ đích: xét nghiệm đường huyết là căn cứ chẩn đoán và do \agent{KnowledgeAgent} đối chiếu với ngưỡng ADA, còn mô hình trả lời câu hỏi khác là ``với hồ sơ này, nguy cơ mắc ĐTĐ cao đến mức nào''. Hai góc nhìn độc lập này là nguyên liệu cho hội đồng kiểm tra chéo.

\subsection{Mô hình và huấn luyện}

Dữ liệu chia 70/10/20 theo phân tầng thành ba phần tách rời: huấn luyện, kiểm định (validation) và kiểm tra (14.451 / 2.065 / 4.130 người; chi tiết cho cả hai bộ dữ liệu ở \tabref{tab:split}). Hai mô hình cùng được huấn luyện:
\begin{itemize}
  \item \textbf{ML}: Gradient Boosting dạng histogram (\texttt{HistGradientBoostingClassifier}) \cite{friedman2001,ke2017}, xử lý giá trị thiếu tự nhiên.
  \item \textbf{DL}: mạng MLP 13 → 128 → 64 → 32 → 1, mỗi lớp ẩn có BatchNorm, ReLU, Dropout 0,2; tối ưu AdamW (lr $10^{-3}$, weight decay $10^{-4}$), batch 512, dừng sớm theo AUC validation (patience 8) \cite{paszke2019,goodfellow2016}.
\end{itemize}
Người dùng thật hiếm khi khai đủ 13 trường. Vì vậy tập huấn luyện được \emph{nhân đôi với bản che ngẫu nhiên}: bản sao bị che 40\% các đặc trưng không bắt buộc (mọi trường trừ tuổi và BMI), để mô hình học cách dự đoán khi thiếu dữ liệu.

\begin{table}[H]
\centering
\caption{Kết quả trên tập kiểm tra NHANES (4.130 người)}
\label{tab:dl-metrics}
\small
\begin{tabular}{@{}lcccc@{}}
\toprule
\textbf{Mô hình} & \textbf{ROC-AUC} & \textbf{PR-AUC} & \textbf{Brier} & \textbf{AUC khi thiếu dữ liệu} \\
\midrule
Hồi quy logistic (đối chứng) & 0,825 & 0,533 & 0,122 & 0,805 \\
ML: Gradient Boosting & 0,817 & 0,523 & 0,124 & 0,802 \\
DL: MLP (PyTorch) & \textbf{0,829} & 0,537 & 0,121 & \textbf{0,809} \\
Trung bình ML + DL (dùng trong agent) & 0,827 & \textbf{0,539} & 0,121 & -- \\
\bottomrule
\end{tabular}
\end{table}

Trên nhóm người gốc Á (561 người trong tập kiểm tra), MLP đạt AUC 0,850, cao hơn toàn bộ tập. Hai mô hình đồng ý về phía của ngưỡng ở 92,3\% ca. Cần nói thẳng: MLP chỉ nhỉnh hơn hồi quy logistic một chút. Với 13 đặc trưng dạng bảng, học sâu không có lợi thế lớn. Lý do vẫn dùng cả hai là để \emph{kiểm tra chéo}: khi ML và DL lệch nhau nhiều, đó thường là dấu hiệu hồ sơ bất thường, và hội đồng sẽ chất vấn (Mục~\ref{sec:council}).

\figw[0.92\linewidth]{ch2_roc_pr}{Đường ROC và Precision–Recall của các mô hình trên tập kiểm tra NHANES}

\subsection{Dữ liệu quy mô lớn: BRFSS 2011–2023}

Để mô hình học từ một quần thể đủ lớn và đa dạng, khối DB được bổ sung \textbf{BRFSS} (Behavioral Risk Factor Surveillance System), khảo sát sức khỏe qua điện thoại lớn nhất thế giới do CDC thực hiện hằng năm. Bản BRFSS 2015 đã làm sạch (253.680 dòng, một năm) là bộ dữ liệu phổ biến cho bài toán dự đoán ĐTĐ. Project không dùng bản đó mà tự xử lý \textbf{dữ liệu gốc của bảy năm lẻ 2011–2023} bằng \texttt{scripts/fetch\_brfss.py}. Chỉ năm lẻ được dùng vì câu hỏi về tăng huyết áp và cholesterol chỉ có trong các năm này. Kết quả là \textbf{3.172.819 người}, gấp 12,5 lần bản phổ biến (\tabref{tab:brfss-years}).

\begin{table}[H]
\centering
\caption{Dữ liệu BRFSS theo năm sau khi chuẩn hóa}
\label{tab:brfss-years}
\small
\begin{tabular}{@{}lrcc@{}}
\toprule
\textbf{Năm} & \textbf{Số người} & \textbf{ĐTĐ} & \textbf{Tiền ĐTĐ} \\
\midrule
2011 & 504.530 & 12,4\% & 1,6\% \\
2013 & 490.942 & 12,7\% & 1,8\% \\
2015 & 440.658 & 13,0\% & 1,7\% \\
2017 & 449.219 & 13,5\% & 1,8\% \\
2019 & 417.423 & 13,8\% & 2,2\% \\
2021 & 437.708 & 13,2\% & 2,3\% \\
2023 & 432.339 & 13,8\% & 2,5\% \\
\midrule
\textbf{Tổng} & \textbf{3.172.819} & 13,2\% & 2,0\% \\
\bottomrule
\end{tabular}
\end{table}

Tên biến của BRFSS thay đổi qua các năm (\texttt{DIABETE2/3/4}, \texttt{\_RFCHOL/1/2/3}, \texttt{SEX/SEXVAR}\ldots). Script dò biến theo danh sách ứng viên cho từng năm và suy ra biến còn thiếu (năm 2011 chưa có \texttt{\_MICHD}, nên bệnh tim được suy từ nhồi máu cơ tim và bệnh mạch vành). Mười một đặc trưng được đặt đúng tên trường của hồ sơ bệnh nhân để agent ánh xạ trực tiếp: tuổi, giới, BMI, tăng huyết áp, rối loạn mỡ máu, hút thuốc, vận động, bệnh tim, đột quỵ, uống nhiều rượu bia, sức khỏe tự đánh giá. Nhãn ĐTĐ là \emph{tự khai} (``bác sĩ từng nói bạn bị tiểu đường'', không tính ĐTĐ thai kỳ); câu trả lời ``không biết'' bị loại. Toàn bộ 3,17 triệu dòng được nạp vào bảng \texttt{brfss\_samples} của SQLite rồi huấn luyện cùng quy trình với NHANES, MLP dùng batch 4.096 và dừng sớm sau 10 epoch. Toàn bộ 3.172.819 dòng đều được dùng, nhưng mỗi phần một vai trò riêng (\tabref{tab:split}).

\begin{table}[H]
\centering
\caption{Chia dữ liệu huấn luyện – kiểm định – kiểm tra (70/10/20, phân tầng theo nhãn, \texttt{random\_state=42})}
\label{tab:split}
\small
\begin{tabularx}{\linewidth}{@{}L{3.3cm}rrY@{}}
\toprule
\textbf{Phần} & \textbf{NHANES} & \textbf{BRFSS} & \textbf{Vai trò} \\
\midrule
Huấn luyện (70\%) & 14.451 & 2.220.973 & Mô hình học trực tiếp \\
\quad sau tăng cường (×2) & 28.902 & 4.441.946 & Thêm bản sao bị che 40\% đặc trưng không bắt buộc, để học cách dự đoán khi người dùng khai thiếu \\
Kiểm định (10\%) & 2.065 & 317.282 & Chọn epoch, giảm tốc độ học, dừng sớm MLP; không dùng để cập nhật trọng số \\
Kiểm tra (20\%) & 4.130 & 634.564 & Chỉ dùng một lần để báo cáo AUC, PR-AUC, Brier (\tabref{tab:dl-metrics}, \tabref{tab:brfss-metrics}) \\
\midrule
\textbf{Tổng} & \textbf{20.646} & \textbf{3.172.819} & \\
\bottomrule
\end{tabularx}
\end{table}

Tập kiểm tra được giữ tách riêng từ đầu đến cuối, nên các chỉ số trong báo cáo phản ánh khả năng dự đoán trên người mà mô hình chưa từng thấy. Tập kiểm tra bị che (``AUC khi thiếu dữ liệu'') được tạo từ chính tập kiểm tra này. Riêng phần đánh giá đầu–cuối của hệ đa agent (Chương~\ref{chap:code}) lấy 400 ca từ tập kiểm tra NHANES.

\begin{table}[H]
\centering
\caption{Kết quả trên tập kiểm tra BRFSS (634.564 người)}
\label{tab:brfss-metrics}
\small
\begin{tabular}{@{}lcccc@{}}
\toprule
\textbf{Mô hình} & \textbf{ROC-AUC} & \textbf{PR-AUC} & \textbf{AUC khi thiếu dữ liệu} & \textbf{AUC người gốc Á} \\
\midrule
Hồi quy logistic (đối chứng) & 0,825 & 0,399 & 0,804 & 0,843 \\
ML: Gradient Boosting & \textbf{0,833} & \textbf{0,422} & \textbf{0,813} & \textbf{0,849} \\
DL: MLP (PyTorch) & \textbf{0,833} & 0,422 & \textbf{0,813} & 0,848 \\
\bottomrule
\end{tabular}
\end{table}

\textbf{Nhận xét.} (i) Trên 634.564 người kiểm tra, cả GBDT và MLP đạt AUC 0,833 và vượt hồi quy logistic rõ hơn so với trên NHANES. Với dữ liệu lớn, mô hình phi tuyến có đủ mẫu để học tương tác giữa các yếu tố. (ii) Nhóm người gốc Á có 14.249 người trong tập kiểm tra, gấp 25 lần NHANES (561), và AUC đạt 0,849. Đây là căn cứ tốt hơn nhiều cho câu hỏi ``mô hình có dùng được cho người châu Á không''. (iii) Dù gấp 150 lần dữ liệu, AUC chỉ ngang NHANES (0,83). Giới hạn nằm ở \emph{thông tin} trong các đặc trưng tự khai, không nằm ở số mẫu. Vì vậy NHANES, với nhãn xác nhận bằng xét nghiệm và số đo thật (vòng eo, huyết áp), vẫn là mô hình chính, còn BRFSS là ý kiến thứ hai.

\subsection{Từ mô hình thành agent}

Một con số xác suất chưa đủ để tư vấn. \agent{PredictionAgent} bọc hai mô hình thành một agent có bốn năng lực:
\begin{enumerate}
  \item \textbf{Dự đoán}: xác suất trung bình của ML và DL, kèm \emph{nguy cơ tương đối} so với tỉ lệ nền 19,8\% của quần thể.
  \item \textbf{Giải thích}: tác động của từng đặc trưng, tính bằng cách thay giá trị đó bằng giá trị tham chiếu rồi đo xác suất thay đổi bao nhiêu.
  \item \textbf{Tự đánh giá}: độ đầy đủ của dữ liệu (tỉ lệ đặc trưng có giá trị thật) và độ lệch giữa hai mô hình.
  \item \textbf{Tranh biện}: trả lời bốn loại \texttt{CHALLENGE} (chất lượng dữ liệu, bất đồng ML–DL, mâu thuẫn với xét nghiệm, bất đồng giữa hai quần thể dữ liệu) bằng cách bảo vệ, nhượng bộ hoặc hạ độ tin cậy.
  \item \textbf{Ý kiến thứ hai quy mô lớn}: chạy thêm cặp GBDT + MLP của BRFSS. Hai bộ dữ liệu có tỷ lệ nền khác nhau (19,8\% và 13,2\%), nên hai mô hình được so theo \emph{nguy cơ tương đối} (gấp bao nhiêu lần trung bình quần thể của mình), không so xác suất tuyệt đối.
\end{enumerate}
Ngoài ra agent chạy mạng hồi quy (RNN) khi người dùng tải chuỗi đo glucose ≥ 5 điểm, để tính các chỉ số theo dõi liên tục (TIR, TBR, CV). Agent không sinh dữ liệu giả khi thiếu đầu vào.

\textbf{Phân loại.} Kết quả dự đoán cùng với xét nghiệm được hội đồng phân tầng thành bốn mức: \emph{Rất cao} (xét nghiệm ở ngưỡng ĐTĐ), \emph{Cao} (tiền ĐTĐ hoặc xác suất ≥ 30\%), \emph{Trung bình} (15–30\%), \emph{Thấp}. Xét nghiệm là căn cứ chính, mô hình bổ sung nguy cơ dài hạn.

\section{Knowledge Base và RAG}
\label{sec:kb}

\subsection{Nguồn tri thức và quy trình tạo KB}

Khối ``web, book → tạo → neo4j KB'' được hiện thực bằng \texttt{scripts/ingest\_kb.py}. Nguồn là 10 tài liệu tiếng Việt biên soạn từ hướng dẫn ADA 2024 \cite{ada2024} và QĐ 5481/QĐ-BYT \cite{byt5481}: tiêu chuẩn chẩn đoán, tiền ĐTĐ, mục tiêu điều trị, sàng lọc người Việt, triệu chứng cảnh báo, biến chứng, thuốc, lối sống và dinh dưỡng, theo dõi glucose liên tục. Có thể nạp thêm bất kỳ trang web hay PDF nào bằng \texttt{--url}. Quy trình:

\begin{figure}[H]
\centering
\begin{tikzpicture}[font=\scriptsize,
  box/.style={draw, rounded corners=2pt, align=center, minimum height=10mm, text width=2.3cm, inner sep=2pt},
  arr/.style={-{Stealth[length=1.6mm]}, black!60}]
\node[box, fill=gray!10] (src) at (0,0) {web, sách, PDF\\Markdown, URL};
\node[box] (split) at (3.1,0) {Trích văn bản\\chia đoạn theo mục};
\node[box, fill=blue!6] (sql) at (6.2,1.0) {SQLite\\\texttt{kb\_chunks} + FTS5};
\node[box, fill=blue!6] (qd) at (6.2,-0.4) {embedding 384 chiều\\→ Qdrant};
\node[box] (tri) at (3.1,-2.0) {Bộ ba tri thức\\tuyển chọn (+ LLM)};
\node[box, fill=orange!10] (neo) at (6.2,-2.0) {MERGE vào Neo4j\\cạnh gắn \texttt{chunk\_ids}};
\node[box, fill=gray!10] (s3) at (9.3,0) {Bản gốc → SeaweedFS};
\draw[arr] (src) -- (split); \draw[arr] (split) -- (sql); \draw[arr] (split) -- (qd);
\draw[arr] (split) -- (tri); \draw[arr] (tri) -- (neo);
\draw[arr, dashed] (src.north) to[bend left=18] (s3.north);
\end{tikzpicture}
\caption{Quy trình tạo Knowledge Base từ web và sách}
\label{fig:ingest}
\end{figure}

Kết quả là 36 đoạn văn và một đồ thị tri thức \textbf{34 nút, 40 cạnh} trên Neo4j. Ontology gồm 6 nhãn nút (\texttt{RiskFactor}, \texttt{Disease}, \texttt{Complication}, \texttt{Test}, \texttt{Drug}, \texttt{Intervention}) và 7 loại quan hệ (\texttt{HAS\_RISK\_FACTOR}, \texttt{DIAGNOSED\_BY}, \texttt{CAUSES}, \texttt{SCREENED\_BY}, \texttt{TREATED\_BY}, \texttt{CONTRAINDICATED\_IN}, \texttt{PREVENTED\_BY}). Mỗi cạnh lưu đoạn trích làm bằng chứng và \texttt{chunk\_ids} trỏ về đoạn văn gốc, nên khi GraphRAG đi qua một cạnh, hệ thống trích dẫn được đúng đoạn tài liệu. Quá trình nạp chạy lại an toàn: tài liệu không đổi (cùng mã băm và cùng mô hình embedding, và đã có trong Qdrant) sẽ được bỏ qua.

\figw[0.85\linewidth]{ch5_kg}{Đồ thị tri thức về ĐTĐ (nút theo nhãn của ontology)}

\subsection{Truy xuất: tìm kiếm lai và GraphRAG}

\agent{KnowledgeAgent} kết hợp ba cách truy xuất \cite{lewis2020,gao2023rag}:
\begin{itemize}
  \item \textbf{Vector} trên Qdrant: bắt các câu hỏi diễn đạt khác từ ngữ trong tài liệu (``khát nước nhiều'' và ``uống nhiều'').
  \item \textbf{Toàn văn} BM25 \cite{robertson2009} trên SQLite FTS5: bắt đúng từ khóa chuyên môn (HbA1c, SGLT2, eGFR). Câu hỏi tự nhiên được nối từ bằng \texttt{OR}, bỏ các từ hư (``bao nhiêu'', ``thì'', ``là''), và tiêu đề mục có trọng số gấp 3 nội dung.
  \item \textbf{GraphRAG} \cite{edge2024graphrag} trên Neo4j: nhận diện thực thể trong hồ sơ hoặc câu hỏi, duyệt lân cận 1–2 bước bằng Cypher, lấy các đoạn văn gắn với cạnh.
\end{itemize}
Hai danh sách vector và toàn văn được gộp bằng Reciprocal Rank Fusion \cite{cormack2009}: $\mathrm{score}(d)=\sum_{r}\frac{1}{60+\mathrm{rank}_r(d)}$. Ngoài truy xuất văn bản, agent còn đối chiếu xét nghiệm của bệnh nhân với ngưỡng ADA (FPG, HbA1c, OGTT, glucose bất kỳ) và kiểm tra quy tắc ``hai kết quả bất thường'' để biết đã đủ tiêu chí chẩn đoán hay cần xét nghiệm lại.

\begin{code}{GraphRAG trên Neo4j (\texttt{src/storage/graph\_store.py}, rút gọn)}
\label{lst:neo4j}
\begin{minted}{python}
def neighborhood(self, start_names, max_hops=2, limit=60):
    records = self._run(f"MATCH p = (s)-[*1..{int(max_hops)}]-(t) WHERE s.name IN $names "
                        f"RETURN p LIMIT $lim", names=list(start_names), lim=limit * 4)
    edges = {}
    for rec in records:
        for hop, rel in enumerate(rec["p"].relationships, 1):   # cạnh thứ hop trên đường đi
            if rel.element_id in edges:
                edges[rel.element_id]["hops"] = min(edges[rel.element_id]["hops"], hop)
                continue
            edges[rel.element_id] = {
                "head": rel.start_node["name"], "relation": rel.type, "tail": rel.end_node["name"],
                "evidence": rel.get("evidence", ""), "chunk_ids": list(rel.get("chunk_ids") or []),
                "hops": hop}
    return sorted(edges.values(), key=lambda e: e["hops"])[:limit]
\end{minted}
\end{code}

\subsection{Đánh giá truy xuất}

Bộ đánh giá gồm 60 câu hỏi tiếng Việt có đáp án (đoạn văn đúng được gắn tay), chia bốn loại: từ khóa (20), diễn đạt lại (19), con số (17), nhiều bước (4). \tabref{tab:rag} đo trên bộ lưu trữ hiện tại của project.

\begin{table}[H]
\centering
\caption{Truy xuất trên 60 câu hỏi tiếng Việt (36 đoạn văn)}
\label{tab:rag}
\small
\begin{tabular}{@{}lcccc@{}}
\toprule
\textbf{Phương pháp} & \textbf{hit@1} & \textbf{hit@3} & \textbf{hit@5} & \textbf{MRR@10} \\
\midrule
Vector (Qdrant, MiniLM) & 0,567 & 0,750 & 0,833 & 0,675 \\
Toàn văn (SQLite FTS5, BM25) & \textbf{0,833} & \textbf{0,967} & \textbf{0,983} & \textbf{0,899} \\
Lai RRF (Qdrant + FTS5, đang dùng) & 0,717 & 0,900 & 0,933 & 0,820 \\
\midrule
\emph{Tham chiếu: lai pgvector + tsvector (bản trước)} & 0,667 & 0,850 & 0,917 & 0,774 \\
\bottomrule
\end{tabular}
\end{table}

\textbf{Nhận xét.} (i) Bộ lưu trữ mới tốt hơn bản PostgreSQL trước đó (hit@3 0,90 so với 0,85), chủ yếu nhờ BM25 của FTS5 cùng xử lý từ hư và trọng số tiêu đề. Độ trễ trung vị của một truy vấn lai là 2,8 ms. (ii) Trên bộ câu hỏi này, toàn văn đơn lẻ lại tốt nhất. Nguyên nhân là mô hình embedding MiniLM yếu với tiếng Việt chuyên ngành (hit@3 0,75), nên khi gộp, nó kéo xếp hạng của BM25 xuống. Phần thí nghiệm của tiểu luận cho thấy mô hình E5 đa ngôn ngữ đạt hit@3 0,98 trên cùng bộ câu hỏi. Đổi embedding hoặc gán trọng số cho RRF là hướng cải thiện trực tiếp. Hệ thống vẫn giữ tìm kiếm lai vì bộ 60 câu nhỏ và có nhiều câu chứa đúng từ khóa của tài liệu, trong khi người dùng thật thường diễn đạt tự do hơn.

\section{Reinforcement Learning và luật}
\label{sec:rl}

\subsection{Agent gợi ý can thiệp bằng Q-learning}

\agent{TreatmentRLAgent} mô hình hóa việc chọn bước can thiệp như một quá trình quyết định Markov (MDP). Trạng thái là tổ hợp 4 mức glucose × 3 mức BMI × 2 mức bệnh đồng mắc, tức 24 trạng thái. Hành động gồm 4 lựa chọn: (0) thay đổi lối sống, (1) metformin, (2) phối hợp thuốc có lợi ích tim–thận, (3) khám chuyên khoa hoặc insulin. Q-learning \cite{watkins1992,sutton2018} cập nhật giá trị hành động theo sai số thời gian:
\begin{equation}
Q(s,a) \leftarrow Q(s,a) + \alpha\big[r + \gamma\max_{a'}Q(s',a') - Q(s,a)\big],\qquad \alpha=0{,}1,\ \gamma=0{,}9,
\label{eq:qlearn}
\end{equation}
với chính sách khám phá $\epsilon$-tham lam ($\epsilon=0{,}15$), huấn luyện 1.500 tập. Hàm thưởng mô phỏng mức cải thiện đường huyết trừ chi phí và độ xâm lấn: với người glucose bình thường, lối sống được thưởng +4, thuốc bị phạt; với người ở ngưỡng ĐTĐ có bệnh đồng mắc, phối hợp tim–thận được thưởng cao nhất; với đường huyết rất cao, chuyển chuyên khoa được +5.

\textbf{Hạn chế cần nói rõ.} Môi trường và hàm thưởng do nhóm tự đặt theo khuyến cáo, không học từ dữ liệu điều trị thật. Q-learning ở đây chủ yếu minh họa cơ chế ``agent học ra quyết định tuần tự'' trong sơ đồ, và chính sách học được không bao giờ được hành động một mình.

\subsection{Luật ràng buộc chính sách học được}

Khối ``rule'' trên sơ đồ là lớp luật đứng trên chính sách RL. Luật luôn thắng khi hai bên khác nhau:
\begin{itemize}
  \item Hội đồng \emph{chưa} xác định ĐTĐ → hành động bị giới hạn ở thay đổi lối sống.
  \item Chẩn đoán cần xác nhận lại (mới có một xét nghiệm bất thường) → hoãn metformin.
  \item Guardrail khẩn được kích hoạt → chuyển khám chuyên khoa.
  \item Chống chỉ định trong đồ thị tri thức (metformin khi eGFR < 30) → luôn được nêu ra.
\end{itemize}
Agent trả về cả hành động RL gốc (\texttt{rl\_action\_id}), hành động sau ràng buộc và cờ \texttt{overridden\_by\_consensus}, nên người xem thấy được khi nào luật đã ghi đè mô hình học.

\begin{table}[H]
\centering
\caption{Safety Guardrail: luật an toàn luôn được ưu tiên}
\label{tab:guardrail}
\small
\begin{tabularx}{\linewidth}{@{}L{2.4cm}L{6.0cm}Y@{}}
\toprule
\textbf{Mức} & \textbf{Điều kiện} & \textbf{Hành động} \\
\midrule
Cấp cứu & Glucose ≥ 600 mg/dL & Đến cơ sở y tế cấp cứu ngay \\
Khẩn & Glucose ≥ 300 mg/dL hoặc HbA1c ≥ 10\% & Khám chuyên khoa Nội tiết sớm; RL chuyển sang hành động 3 \\
Khẩn & Chuỗi đo có glucose < 54 mg/dL & Cảnh báo hạ đường huyết mức 2 \\
Ưu tiên & BMI ≥ 25 + tăng huyết áp + tiền sử gia đình hoặc tim mạch & Nhóm nguy cơ ``đỏ'': ưu tiên xét nghiệm và kiểm soát yếu tố tim mạch \\
\bottomrule
\end{tabularx}
\end{table}

\section{Hội đồng tranh biện: tư vấn và trợ giúp}
\label{sec:council}

Vòng tròn lớn trên sơ đồ là \agent{ConsensusCriticAgent}. Nó nhận kết quả của hai nhánh, nhưng không chỉ tổng hợp: khi phát hiện mâu thuẫn, nó gửi \texttt{CHALLENGE} qua bus tới agent liên quan, nhận \texttt{RESPONSE} và cập nhật kết luận. Phản hồi ở vòng trước có thể mở ra câu hỏi ở vòng sau, tối đa 2 vòng.

\begin{table}[H]
\centering
\caption{Giao thức tranh biện của hội đồng}
\label{tab:debate}
\small
\begin{tabularx}{\linewidth}{@{}L{3.6cm}L{4.0cm}Y@{}}
\toprule
\textbf{Vấn đề phát hiện} & \textbf{Chất vấn} & \textbf{Phản hồi và ảnh hưởng} \\
\midrule
Dữ liệu thật < 35\% số đặc trưng & Prediction · \texttt{CHALLENGE\_DATA\_QUALITY} & Hạ độ tin cậy, liệt kê đặc trưng còn thiếu → (−) \\
ML và DL lệch ≥ 15 điểm \% & Prediction · \texttt{CHALLENGE\_MODEL\_DISAGREEMENT} & Chỉ ra đặc trưng hai mô hình đánh giá khác nhau nhất → (−) \\
ML ≥ 30\% nhưng xét nghiệm bình thường & Prediction · \texttt{CHALLENGE\_ML\_VS\_LABS} & Giữ quan điểm ``nguy cơ dài hạn'' → (−) và \textbf{mở vòng 2} \\
(vòng 2) ML giữ quan điểm & Knowledge · \texttt{CHALLENGE\_SCREENING\_PLAN} & Lịch tầm soát ADA có trích dẫn → (+) \\
ML < 15\% nhưng xét nghiệm ở ngưỡng ĐTĐ & Prediction · \texttt{CHALLENGE\_ML\_VS\_LABS} & Nhượng bộ cho xét nghiệm → (−) \\
Xét nghiệm ở ngưỡng ĐTĐ & Knowledge · \texttt{CHALLENGE\_VERIFY\_CRITERIA} & Kiểm tra quy tắc ``2 kết quả bất thường'' → ``đủ tiêu chí'' hoặc ``cần xét nghiệm lại'' \\
Mô hình NHANES và BRFSS ở hai phía ngược nhau (một bên ≥ 1,5 lần, bên kia ≤ 0,8 lần trung bình) & Prediction · \texttt{CHALLENGE\_COHORT\_DISAGREEMENT} & Giải trình khác biệt về nhãn và đặc trưng của hai quần thể, giữ mô hình chính, hạ độ tin cậy → (−). Cùng phía thì ghi (+) ``mô hình quy mô lớn đồng thuận'' \\
\bottomrule
\end{tabularx}
\end{table}

\begin{code}{Vòng tranh biện của hội đồng (\texttt{src/agents/consensus\_critic\_agent.py}, rút gọn)}
\label{lst:critic}
\begin{minted}{python}
for round_no in range(1, settings.max_debate_rounds + 1):
    issues = [i for i in _detect_issues(state) if i[0] not in raised]
    if not issues:
        break                                            # không còn mâu thuẫn → kết thúc
    raised.update(i[0] for i in issues)
    replies = await asyncio.gather(*(                    # chất vấn song song qua MessageBus
        self.ask(receiver, msg_type, payload, msg.conversation_id, performative="CHALLENGE")
        for _, receiver, msg_type, payload in issues))
    for (issue_id, receiver, msg_type, _), reply in zip(issues, replies):
        debate.append({"round": round_no, "issue": issue_id, "challenged": receiver, ...})
        if reply.performative == "RESPONSE":
            _apply_response(state, issue_id, reply.payload)   # cập nhật (+)/(−), độ tin cậy
return self.reply(msg, "CONSENSUS_RESULT", _conclude(state, profile, debate))
\end{minted}
\end{code}

Kết luận của hội đồng gồm mức nguy cơ, danh sách bằng chứng đồng thuận (+), điểm mâu thuẫn hoặc cần lưu ý (−), cảnh báo an toàn, và kế hoạch tầm soát nếu có. \agent{ResponseAgent} dùng đúng các dữ liệu đó để viết câu trả lời: phần \emph{tư vấn} (mức nguy cơ và lý do, khuyến cáo có trích dẫn nguồn) và phần \emph{trợ giúp} (bước tiếp theo nên làm, câu hỏi làm rõ khi thiếu dữ liệu). Agent bị cấm thêm thông tin không có trong kết quả của các agent khác, và luôn nhắc rằng hệ thống không thay thế bác sĩ.

\section{Tiếp nhận, làm rõ và bộ nhớ}

Ba agent ở tầng tiếp nhận biến dữ liệu thô thành hồ sơ có cấu trúc:
\begin{itemize}
  \item \agent{ConversationAgent} bóc tách chỉ số từ câu tự do, ví dụ ``mẹ tôi 61 tuổi, nặng 68 ký cao 1m52, sáng nay thử đường máu lúc chưa ăn được 6,4''. Đơn vị không ghi rõ được suy ra (6,4 → mmol/L → 115 mg/dL) kèm cờ \texttt{unit\_inferred}.
  \item \agent{DocumentAgent} đọc phiếu xét nghiệm dạng PDF, ảnh chụp, DOCX hoặc CSV chuỗi đo glucose.
  \item \agent{ClarificationAgent} phát hiện hồ sơ ngoài phạm vi sàng lọc (thai phụ, người đã được chẩn đoán, đang dùng thuốc hạ đường huyết, HbA1c lệch nhiều so với glucose\ldots) và đặt tối đa 2 câu hỏi làm rõ mỗi lượt, không hỏi lại câu đã hỏi.
\end{itemize}
Mỗi trường trong \texttt{PatientProfile} lưu giá trị, nguồn (\texttt{text}, \texttt{file:<tên>}), độ tin cậy, đoạn trích làm bằng chứng và cờ cảnh báo. Hệ thống không tự điền giá trị mặc định cho chỉ số lâm sàng. Khi hai nguồn mâu thuẫn, hội đồng báo (−) để người dùng xác nhận. Bộ nhớ hai tầng (hồ sơ tích lũy và 10 tin nhắn gần nhất) giúp hệ thống hiểu câu trả lời cụt (``62'' sau khi được hỏi tuổi) và trả lời câu hỏi tiếp nối mà không chạy lại hội đồng.
